\emph{Frames.} Near a point of $\Sig$ at which $G$ has corank one, choose a smooth frame $(k, m_1, \dots, m_{n-1})$ of $T^*M$ whose first member spans the radical on $\Sig$. The block $G_1$ of the $m_i$ is nondegenerate on $\Sig$ (its kernel would lie in the radical, which meets $\operatorname{span}(m_i)$ only in $0$), hence near $\Sig$, and the block formula gives $\det G = \det G_1 \cdot S$ with $S = \lambda - b^{\mathsf T} G_1^{-1} b$. Since $k$ is radical on $\Sig$, $\lambda$ and $b$ vanish there. A change of frame multiplies $\det G$ by the square of a nowhere zero function, so the order of vanishing of $\det G$ along $\Sig$, and whether its zero is regular, are frame independent; by $\det G = \det G_1 S$ the same holds for $S$. The inertia satisfies $\operatorname{In}(G) = \operatorname{In}(G_1) + \operatorname{sgn} S$ off $\Sig$ (congruence by the block elimination, inertia being invariant under congruence by \claimref{synthesis.completion-recovers-the-real-clifford-algebras}(iii)).

\emph{(1)} Jacobi's formula, $d \det G = \operatorname{tr}(\operatorname{adj}(G)\, dG)$, holds for any smooth family of square matrices (differentiate the determinant column by column and expand along the differentiated column). At a point of corank one in the frame above the Gram matrix is $0 \oplus G_1$, whose adjugate is $\det G_1 \cdot E_{00}$ ($E_{00}$ the matrix unit of the first slot): every $(n{-}1)$-minor that omits the first row or column contains the zero row or column, except the complementary minor of the $(0,0)$ entry. So $d\det G = \det G_1 \cdot (dG)(k,k)$. Extending $k$ as in Lemma \claimref{sheaf.invariant-form-t3}, $d(G(k,k)) = (dG)(k,k) + 2G(\nabla k, k)$ in any trivialisation, and the second term vanishes on $\Sig$ since $k$ is radical there; so $(dG)(k,k)|_\Sig = dg_k|_\Sig = \delta(k\otimes k)$ and $d\det G = \det G_1\,\delta(k\otimes k)$. At a point of corank $r \geq 2$ every $(n{-}1)$-minor vanishes (rank at most $n - 2$), so $\operatorname{adj}(G) = 0$ and $d \det G = 0$.

Now suppose (T1)-(T3). Off $\Sig$ the form is nondegenerate, so the zero set of $\det G$ is $\Sig$; on $\Sig$ the corank is one (T2) and $\delta \neq 0$ (T3), so $d\det G \neq 0$. Conversely suppose $\det G$ vanishes exactly on $\Sig$ with nonzero differential there. Then $\Sig$ is a regular level set, hence a smooth hypersurface, closed as the zero set of a continuous function: (T1). At each of its points the corank is one, because corank at least two forces $d\det G = 0$; the radical line is the column space of $\operatorname{adj}(G)$, which has rank exactly one at corank one ($G \operatorname{adj}(G) = \det G \cdot 1 = 0$ puts its columns in the kernel, and the minor complementary to a nonzero entry of the radical is nonzero), so a nonvanishing column is a smooth local spanning section: (T2). Finally $\delta(k\otimes k) = d\det G/\det G_1 \neq 0$: (T3). Since $\det G = \det G_1 S$ with $\det G_1 \neq 0$, (T3) is the statement that $S$ has a simple zero, $\partial_\tau S \neq 0$ on $\Sig$.

\emph{(2)} Take collar data $(k, u, \beta, \mu, \nu)$ with $\mu\nu > 0$ (Definition \claimref{obj.sheaf.collar-data-coupling}; they exist since $w_1(\shL) = 0$, and the definition uses only (T1), (T2)), and work in the frame above along transversals, so that $e$ has components $(1, e')$ in the dual frame ($k(u) = 1$). The block inverse of $\begin{psmallmatrix} \lambda & b^{\mathsf T} \\ b & G_1\end{psmallmatrix}$ is
\[
\begin{pmatrix} S^{-1} & -S^{-1} b^{\mathsf T} G_1^{-1} \\ -S^{-1} G_1^{-1} b & G_1^{-1} + S^{-1} G_1^{-1} b\, b^{\mathsf T} G_1^{-1} \end{pmatrix},
\]
so $s = e^{\mathsf T} G^{-1} e = S^{-1}(1 - 2c + c^2) + B = (1-c)^2/S + B$ exactly, with $c$ and $B$ smooth up to $\Sig$ (built from $b$, $G_1^{-1}$ and $e'$) and $c|_\Sig = 0$ because $b|_\Sig = 0$. By compactness of $\Sig$ there is a collar on which $(1-c)^2 \geq \tfrac14$ and $|B| \leq C$, and since $S \to 0$ uniformly at $\Sig$ with $S > 0$ on the physical side, a thinner collar on which $1/(4S) \geq 2C$ there; on it $s \geq 1/(8S) > 0$. Lemma \claimref{sheaf.determinant-identity-blow-up-expansion}(1), an exact Schur complement identity valid wherever $G$ is invertible, gives $\det h = \det G \cdot (-(1 + 2\mu\nu\beta^2 s))$ with $1 + 2\mu\nu\beta^2 s \geq 1$: $h$ is nondegenerate on the physical half of the collar. Outside the collar $h = G \perp H$, nondegenerate of signature $(p{+}1, q{+}1)$. On $\Sig$, put $V_1 := \{\alpha : \alpha(u) = 0\}$, a complement of $\shL$; $h|_{V_1} = G|_{V_1}$ is nondegenerate (it is the induced form on $T^*M/\shL$), of signature $(p{-}1,q)$ because $\operatorname{In}(G_1) = \operatorname{In}(G) - \operatorname{sgn} S = (p{-}1,q)$ on the physical side and $G_1$ stays nondegenerate up to $\Sig$. The coupling pairs $V_1$ with nothing ($\alpha(u) = 0$) and $G(V_1, k) = 0$ on $\Sig$, so $T = V_1 \oplus \operatorname{span}(k, \nO, \nInf)$ $h$-orthogonally, and the Gram matrix of $(k, \nO, \nInf)$ is $\begin{psmallmatrix} 0 & \nu & \mu \\ \nu & 0 & -1 \\ \mu & -1 & 0\end{psmallmatrix}$, of determinant $-2\mu\nu \neq 0$, containing the hyperbolic plane $\operatorname{span}(k, \nO)$ and the vector $k + \nu\nInf - \mu\nO$ of square $2\mu\nu > 0$: signature $(2,1)$. So $h|_\Sig$ is nondegenerate of signature $(p{+}1, q{+}1)$. The signature is constant on $\Mp$ by continuity along the transversals, as in the original proof. Nowhere was the order of vanishing of $S$ used: the expansion $s = 1/\lambda + O(1)$ of the original argument, which is where (T3) entered, is replaced by the exact identity.

\emph{(3)} Let $S = \tau^m u$, $m$ finite, $u$ nowhere zero. Write $N := (1-c)^2$. On the far side, for odd $m$, $S < 0$ and $s = N/S + B \to -\infty$ uniformly. Differentiating along the transversal,
\[
s' \;=\; -\frac{N S'}{S^2} + \frac{N'}{S} + B', \qquad S' = \tau^{m-1}(m u + \tau u'),
\]
so $|N S'/S^2| \geq \tfrac14 (m|u| - |\tau u'|)/(u^2 |\tau|^{m+1})$, which dominates $|N'/S| = O(|\tau|^{-m})$ and $|B'| = O(1)$ as $\tau \to 0$, uniformly by compactness. This is the monotone zone of the induction argument, obtained from finite order instead of (T3). From here the argument runs as for walls satisfying (T3): with $\sigma := \sup |s|$ over the far-side part of the collar outside the monotone zone (finite by compactness) and a good coupling $|\mu\nu| < 1/(2\sigma)$, $F := 1 + 2\mu\nu\beta^2 s$ is positive outside the monotone zone, tends to $-\infty$ at the wall, and inside the zone (where $\beta \equiv 1$ after shrinking) is strictly monotone along each transversal with $F' = 2\mu\nu s' \neq 0$. Its zero set $\Sig'$ is a regular level set, a smooth closed hypersurface, and $\det h = \det G \cdot (-F)$ has a simple zero across it ($\det G \neq 0$ off $\Sig$). By (1), applied to the form $h$ on the bundle $T$ (the argument of (1) uses only a symmetric form and its radical, not the cotangent origin of the bundle), $(T, h)$ satisfies (T1)-(T3) along $\Sig'$. For even $m$, $S > 0$ on both sides of $\Sig$, so the argument of (2) applies on both sides and $h$ is nondegenerate on a neighbourhood of $\Sig$; the signature of $G$ is the same on both sides ($\operatorname{In}(G) = \operatorname{In}(G_1) + \operatorname{sgn} S$).

\emph{(4)} If the signature of $G$ changes across $\Sig$, then at each $x \in \Sig$ the two local sides of $\Sig$ in a small ball carry different signatures, so the normal direction pointing into the side of signature $(p,q)$ is well defined and locally constant: a global co-orientation, so $N\Sig$ is trivial and $\Sig$ is two-sided. This uses neither (T3) nor the canonical isomorphism $N^*\Sig \cong \shL^{\otimes 2}$. At even order a one-sided wall exists: on the open M\"obius band with fibre coordinate $\tau$ and core circle $\Sig = \{\tau = 0\}$, the form $G = \tau^2\, \partial_\tau \otimes \partial_\tau + \partial_x \otimes \partial_x$ is well defined (the identification $\tau \mapsto -\tau$ fixes $\tau^2$ and $\partial_\tau \otimes \partial_\tau$), has corank exactly one on $\Sig$ with smooth radical line, and $S = \tau^2$.

\emph{Final sentence.} For the wild wall with $S = \tau^3$, (2) and (3) apply with $m = 3$, while $\lim_{\tau \to 0} \tau\, G^{-1}$ does not exist: the top-left entry of $G^{-1}$ is $1/S = \tau^{-3}/u$. The check runs this model exactly: $G_1$ negative definite of rank three, $b = \tau b_1$, $\lambda = \tau^3 + \tau^2 b_1^{\mathsf T} G_1^{-1} b_1$, so that $S = \tau^3$, $G(k,k)$ has no linear term and $\det G = \det G_1\, \tau^3$; the identity of (2) at rational points; signature $(2,4)$ of $h$ at the wall and on the physical side; $\det h$ a cubic in $\tau$ with no repeated factor ($\gcd(P, P') = 1$), changing sign between $\tau = -1/4$ and $\tau = -1/5$, with signature $(2,4)$ between the wall and that zero and $(1,5)$ beyond it.

\emph{Which consumers need which clause} (a reading of their stated proofs, not a theorem). The one-sided completion uses (T3) only through the expansion replaced in (2). The induction to shadow walls uses it only for the monotone zone, replaced in (3) by finite order. Two-sidedness follows from the signature change, (4). The absolute form of conservation of degeneracy uses only that the two sides carry different signatures. The degeneration sheaf uses (T1) and (T2). The canonical simple pole of $G^{-1}$, hence the identification $N^*\Sig \cong \shL^{\otimes 2}$ through the residue, and the exceptional algebra of the normal-cone deformation (whose rescaling $k \mapsto k/\sqrt t$ needs $\dot G(k,k) \neq 0$) do use first order, as do the results that import that pole.
